% Pour toute question sur l'utilisation de cette classe, écrivez à jerome.soucy(at)ulaval.ca
% Nécessite la classe Evaluation.sty dans le même dossier ou dans le PATH de votre système
% Nécessite le fichier logo_ul.pdf dans le même dossier un dans le PATH de votre système
\def\Ponderation{10}
\def\SigleNumero{STT-2100}
\def\NomCours{Régression}
\def\DateEvaluation{2 mars 2023}
\def\DateRemise{6 décembre 2023}
\def\HeureRemise{23h59}
\def\HeureEvaluation{}
\def\Session{Été 2021}
\def\NomEvaluation{Devoir 2}

% Ne rien mettre entre les accolades si l'enseignant (ou l'enseignante) est aussi le coordonnateur (ou la coordonnatrice)
\def\Coordonnatrice{}
\def\Coordonnateur{}

% Ne rien mettre entre les accolades s'il y a plus d'une section
\def\Enseignant{Julien Miron}
\def\Enseignante{}

% Ne rien mettre entre les accolades s'il s'agit d'un cours à section unique
% Mettre le nom de chacun des enseignants et enseignantes entre les accolades
% Une case à cocher sera générée pour chacune des sections
\def\SectionA{}
\def\SectionB{}
\def\SectionC{}
\def\SectionD{}


\def\Directives{
\begin{enumerate}
\item Le devoir compte \numquestions ~exercices pour un total de \numpoints{}~points et est à faire individuellement ou en équipe de deux personnes. 

\item Vous devez remettre \underline{deux fichiers} par équipe:
\begin{itemize}
\item Un fichier PDF avec vos réponses aux questions 1 à 3. Vous pouvez rédiger à l'ordinateur dans \LaTeX, à la main sur un papier que vous numérisez ou sur une tablette. Ce fichier est à remettre sur Gradescope, à l'aide des instructions données au dernier point de cette page.
\item Un script \texttt{R} avec les codes utilisés pour la question \ref{prog}. À remettre dans la boîte de dépôt du site du cours.
\end{itemize}

\item Il est \textit{\textbf{nécessaire}} de justifier toutes les démarches. Une réponse qui n'a pas une démarche \textit{\textbf{complète}} ne pourra avoir tous les points. Pour certaines parties du devoir, une démarche peut être une justification en mots.
\item La qualité de la langue sera évaluée, incluant l'orthographe, la grammaire et la syntaxe. Un maximum de 10 points pourra être retranché pour mauvaise qualité du français. 
\item Les travaux remis en retard recevront la note 0.

\item Remises sur \textit{Gradescope}:
\begin{itemize}
\item Rendez-vous au www.gradescope.com;
\item Connectez-vous en utilisant votre adresse courriel de l'Université (idul@ulaval.ca) ainsi que votre mot de passe habituel. Il faut faire «Log in», utiliser «School credentials» et chercher Université Laval dans la liste . Vous devriez déjà avoir un compte, par le biais de l’université. 
\item Une fois sur le site, choisissez STT-2100;
\item En sélectionnant le Devoir 2, vous ajoutez le PDF qui contient vos réponses;
\item Vous devez identifier les différentes sous-questions sur les pages de votre devoir.
\end{itemize}
\end{enumerate}
}

% Veuillez indiquer le type d'évaluation entre crochets
% Examen : option à choisir pour un examen présentiel qui sera corrigé de manière classique. Un tableau des points sera affiché sur la page couverture.
% Gradescope : option à choisir pour un examen qui sera corrigé via Gradescope. Aucun tableau de pointage n'apparaîtra et il y a quelques différences au niveau de la mise en page
% TPS :
% TPF :
% Devoir :
\documentclass[Devoir]{Evaluation}
\usepackage{multicol,graphicx}
\begin{document}

\begin{questions}  
\question Soit le modèle de régression linéaire multiple 
\begin{equation*}\label{oldmod}
    \boldsymbol{Y}=\boldsymbol{X\beta}+\boldsymbol{\varepsilon},
\end{equation*}
où $\boldsymbol{\varepsilon}=(\varepsilon_1,\hdots,\varepsilon_n)^T$, avec $\varepsilon_j\sim \mathcal{U}(-1;1)$, pour $j=1,\hdots,n$, et $\mathcal{U}(a;b)$ est la loi uniforme continue sur l'intervalle $[a;b]$. Nous supposons aussi que $\varepsilon_i\perp \varepsilon_j$, pour tout $(i,j)\in \{1,\hdots,n\}^2$ et $i\neq j$.

\begin{parts}
    \part[2] Calculez l'espérance de $\boldsymbol{Y}$.
    \part[3] Calculez la variance de $\boldsymbol{Y}$.
    \part[2] Donnez la loi de $Y_j$, un des éléments du vecteur $\boldsymbol{Y}$.
    \part[8] Les quatre postulats sont-ils respectés par ce modèle? Expliquez brièvement pour chacun des postulats pourquoi vous dites qu'il est respecté ou non.
    \part[3] Calculez l'estimateur des moindres carrés $\widehat{\boldsymbol{\beta}}_{MC}$ pour $\boldsymbol{\beta}$.
    \part[5] À quel problème sommes-nous confrontés si nous tentons de trouver l'estimateur du maximum de vraisemblance $\widehat{\boldsymbol{\beta}}_{MV}$?
\end{parts}

\vspace{1cm}

\question 
Soit le modèle de régression linéaire multiple 
\begin{equation}\label{oldmod}
    \boldsymbol{Y}=\boldsymbol{X\beta}+\boldsymbol{\varepsilon},
\end{equation}
où $=\boldsymbol{\varepsilon}\sim\mathcal{N}_n(\boldsymbol{0},\sigma^2 \mathbbm{1}_n)$, $\sigma>0$ et $\mathbbm{1}_n$ est la
matrice identité de dimension $n\times n$. 
Pour les besoins de l'exercice, nous écrivons $\boldsymbol{X}=(\boldsymbol{x}_1^T,\hdots,\boldsymbol{x}_n^T)^T$, où $\boldsymbol{x}_i=(1, x_{i1},\hdots, x_{ip})^T$, pour $i=1\hdots,n$. 
Nous aurons donc que $\boldsymbol{\beta}\in \mathbbm{R}^{p+1}$.

En supposant que certaines observations ont des valeurs explicatives «aberrantes», il pourrait être intéressant d'en réduire l'impact dans l'estimation.
Des poids individuels $w_i$ seront donc appliqués aux $\boldsymbol{x}_i$ de façon à avoir $\boldsymbol{X}^*=(w_1\boldsymbol{x}_1^T,\hdots,w_n\boldsymbol{x}_n^T)^T$, où $w_i\boldsymbol{x}_i=(w_i, w_ix_{i1},\hdots, w_ix_{ip})^T$ et $w_i \in [0,1]$, pour $i=1\hdots,n$.  
Le nouveau modèle sera 
\begin{equation}\label{newmod}
    \boldsymbol{Y}=\boldsymbol{X}^* \boldsymbol{\beta}+\boldsymbol{\varepsilon},
\end{equation}
où nous supposons que les $\boldsymbol{\varepsilon}$ sont les mêmes que dans le modèle \eqref{oldmod}.
\begin{parts}
    \part[5] Réécrivez le modèle \eqref{newmod} en fonction de $\boldsymbol{X}$ et identifiez bien toutes les matrices et/ou tous les vecteurs qui sont différents du modèle \eqref{oldmod}.
    \part[3] À l'aide du modèle \eqref{newmod}, nous utilisons $\widehat{\boldsymbol{\beta}}^*=\boldsymbol{(X^{*T}X^*)^{-1}X^{*T}Y}$ comme estimateur de $\boldsymbol{\beta}$. Réécrivez cet estimateur en utilisant la notation que vous avez établie en (a).
    \part[6] Calculez le biais de $\widehat{\boldsymbol{\beta}}^*$.

    \part[7] Calculez la variance de $\widehat{\boldsymbol{\beta}}^*$ et comparez-la avec la variance de $\widehat{\boldsymbol{\beta}}$, l'estimateur des moindres carrés ordinaire trouvé avec le modèle \eqref{oldmod}. Cette comparaison semble-t-elle logique? Expliquez dans vos mots. \textit{Pour comparer, vous pouvez étudier les variances des estimateurs d'un seul coefficient de $\boldsymbol{\beta}$, puis généraliser}.

\end{parts}
\newpage

\question\label{prog}
\textbf{Cet exercice est à compléter sur \texttt{R}. }
Dans le jeu de données fourni pour ce devoir, vous retrouvez les variables $\boldsymbol{Y}$, $\boldsymbol{r}$, $\boldsymbol{s}$ et $\boldsymbol{z}$. Le modèle étudié sera 
\begin{equation}\label{totmod}
    \boldsymbol{Y}=\boldsymbol{X\beta}+\boldsymbol{\varepsilon},
\end{equation}
où $\boldsymbol{X}=(\boldsymbol{1}, \boldsymbol{r}, \boldsymbol{s}, \boldsymbol{z})$ et $\boldsymbol{\varepsilon}\sim \mathcal{N}_{n}(\boldsymbol{0}, \sigma^2 \mathbbm{1}_n)$. Les composants des vecteurs de variables explicatives sont toutes des réalisations de la loi normale centrée-réduite et sont toutes indépendantes. Nous avons aussi que les trois variables explicatives sont indépendantes.

\vspace{0.5cm}

\begin{parts}
    \part[4] Si vous ajustez le modèle \eqref{totmod}, à quelle variable (en excluant l'ordonnée à l'origine) correspond la valeur-p la plus élevée parmi les tests-$t$ sur les coefficients? À noter, lorsqu'on utilise \texttt{summary(modele)\$coefficients}, on obtient les valeur-p exactes au lieu de \texttt{<2e-16}, le cas échéant. Donnez la variable et la valeur-p associée.
    \part[6] Si on vous dit que les données ont bel et bien été générées avec le modèle $$Y_i=\beta_0 + \beta_rr+\beta_ss+\beta_zz +\varepsilon_i,$$ pour $i=1,\hdots,n$, qu'est-ce qui pourrait expliquer que certaines valeurs-p sont supérieures aux seuils standards (e.g. $5\%$)?
    \part[5] Donnez le $VIF$ pour chacune des trois variables explicatives. Que pouvez-vous dire à propos des ces $VIF$? Ces résultats semblent-ils cohérents?


\part[5] À partir d'ici, le modèle étudié sera celui en \eqref{newmod}.
    Il faudra donc définir un poids $w_i$ pour chaque observation du jeu de données.
    Le poids utilisé sera trouvé à partir de la distance de Mahalanobis expliquée dans l'encadré ci-bas.
    Pour ce faire, nous devons utiliser la matrice de design $\boldsymbol{X}$, mais en retirant la colonne de 1. Que vaut la distance de Mahalanobis $d_i$ des cinq premières observations du jeu de données?

    \part[6] Calculez la moyenne des $d_i^2$. Ce résultat vous semble-t-il cohérent avec la théorie? Expliquez.
    \part[10] Les $d_i$ seront maintenant utilisés pour calculer les poids $w_i$ du modèle \eqref{newmod}. 
    Vous devez proposer une pondération qui est fonction de ces $d_i$. 
    Cette pondération n'a pas besoin d'être compliquée, mais on repérera facilement ceux qui utilisent ChatGPT. Donnez la fonction utilisée et expliquez bien comment vous calculez les $w_i$.
    
    \part[10] Implémentez votre poids prososé en (e) et ajustez le modèle \eqref{newmod} en utilisant ces poids. Comparez les résultats avec ceux du modèle en (a) et faites un lien avec le (d) de l'exercice 2.
    
\end{parts}
\vspace{1.5cm}
\fbox{
\parbox{\textwidth}{
\textbf{Distance de Mahalanobis}

\vspace{0.5cm}
Soit $\boldsymbol{u}_i=(u_{i1},\hdots,u_{ik})^T $, pour $i=1,\hdots,n$.
La distance de Mahalanobis $d_i$ de l'observation $i$ est calculée avec
    \begin{equation*}
        d_i=\sqrt{(\boldsymbol{u}_i - \overline{\boldsymbol{u}})^T\Sigma (\boldsymbol{u}_i - \overline{\boldsymbol{u}})}, \text{ pour }i=1,\hdots,n,
    \end{equation*}
où $\overline{\boldsymbol{u}}=(\overline{u}_{1},\hdots,\overline{u}_{k})^T$, avec $\overline{u}_{j}=\frac{1}{n}\sum_{i=1}^n u_{ij}$, pour $j=1,\hdots,k$, et $\Sigma$ est la matrice de covariance empirique des $\boldsymbol{u}_i$. 
Nous savons que si les $\boldsymbol{u}_i$ sont des réalisation de variables aléatoires normales, les statistiques $d_i^2$ sont des réalisations d'une variable aléatoire qui suit la loi du khi-carré avec $k$ degrés de liberté, où $k$ est la longueur des vecteurs $\boldsymbol{u}_1,\hdots, \boldsymbol{u}_n$.
Bien sûr, en \texttt{R}, vous pouvez utiliser la fonction \texttt{mahalanobis()} en utilisant les bons arguments. Il faut toutefois faire attention à la valeur qui est retournée par cette fonction.
}
}
\end{questions}


\end{document}




\vspace{0.3in}
\textbf{Question 4 (10 pts)} (Design d'expérience)\\

On s'intéresse à un médicament sensé réduire le nombre de crises chez les individus atteints d'epilepsie. Vous devez aider à planifier une expérience pour étudier l'effet de la dose du médicament sur le nombre de crises. Puisque vous vous intéressez à la relation entre deux variables quantitatives, vous planifiez analyser les données avec un modèle de régression linéaire simple. \\

Soit $x_1, \ldots, x_n$ les valeurs de la dose pour les 20 individus de votre étude.\\ Pour chacun, vous pouvez attribuer n'importe quelle dose entre 2 mg et 21 mg. \\

\begin{itemize}
%\item [(a)] Pour le modèle de régression linéaire simple $Y_i = \beta_0 + \beta_1 x_i + \varepsilon_i$, prouvez que la variance de $\beta_0$ est minimale si les $x_i$ sont tels que $\bar{x}$ = 0.
\item [(a)] (5 pts) Quelle est la meilleure fa\c{c}on de choisir les valeurs $x_1,\ldots,x_n$ pour minimiser la variance de votre estimateur $\widehat{\beta_1}$? Justifiez votre réponse.
\vspace{0.1 in}
\item [(b)] (5 pts) Pourquoi ne serait-ce pas vraiment une bonne idée de choisir les valeurs de $x_i$ de cette façon, même si cela minimise la variance de votre estimateur?
\end{itemize}

\end{document} 